% !TeX program = xelatex
\documentclass{article}
\usepackage{geometry}
\usepackage[11pt]{moresize}
\usepackage{tikz}

\geometry{paperwidth = 28.875cm, paperheight = 50.8cm, margin = 0cm, inner = 0cm, outer = 0cm}

\usepackage{xfrac}
\usepackage{ifthen}
\usepackage{intcalc}
\usepackage[MnSymbol]{mathspec}	
\setmainfont[Mapping=tex-text,Ligatures={NoRequired,NoCommon,NoContextual}]{Comic Sans MS}
\setmathfont(Digits,Latin,Greek)[Numbers={Lining,Proportional}]{Comic Sans MS}
\usetikzlibrary{shapes.geometric,decorations,decorations.pathmorphing}
\usetikzlibrary{intersections}
\usetikzlibrary{patterns}
\usetikzlibrary{calc}
\usetikzlibrary{shapes.misc}
\usetikzlibrary{spy}

\pgfdeclarelayer{bg}
\pgfdeclarelayer{bbg}
\pgfsetlayers{bbg, bg, main}

\definecolor{c1}{HTML}{ac0000}
\definecolor{c2}{HTML}{00a0d5}

\newcommand{\abs}[1]{{\left|{#1}\right|}}

\tikzset{
	point/.style n args = {1}{circle, draw = black, inner sep = #1pt, fill = black},
	empty point/.style = {circle, draw = black, inner sep = 2pt, fill = white},
	cross/.style = {cross out, draw, minimum size = 2 * (#1 - \pgflinewidth), inner sep = 0pt, outer sep = 0pt},
	xkcd/.style n args = {1}{decorate, decoration = {random steps, pre length = #1mm, post length = #1mm, segment length = 0.6mm, amplitude = 0.2pt}},
	xkcd/.default = {4},
	point/.default = {2}
}

\begin{document}
	\pagestyle{empty}\noindent
	\begin{tikzpicture}
		\begin{pgfonlayer}{bg}
			\draw[opacity = 0.3, dashed, c2] (-14.4375,-25.39) grid (14.4375,25.39);
			\draw[opacity = 0.1, densely dashed, c2, step = 0.5] (-14.4375,-25.39) grid (14.4375,25.39);
		\end{pgfonlayer}
		\node[black] (title) at (0,22) {\HUGE \textbf{Συγκλίνει;}};
		%
		%
		\node[scale = 1.1] (function) at (-7,16) {\begin{tikzpicture}
			\draw[xkcd, thick, ->] (-4,0) -- (4,0) node[pos=1, below]{$x$};
			\draw[xkcd, thick, ->] (0,-4) -- (0,4) node[pos=1, left]{$y$};
			\draw[xkcd, thick, c1] (-4,-3) -- (3,4) node[pos = 1, above] {\scriptsize$f_1$};
			\draw[xkcd, thick, opacity = 0.9, c1] (-4,-3.4) -- (3.4,4) node[pos = 1, above] {\scriptsize$f_2$};
			\draw[xkcd, opacity = 0.7, c1] (-4,-3.7) -- (3.7,4) node[pos = 1, above] {\scriptsize$f_3$};
			\draw[xkcd, opacity = 0.5, c1] (-4,-3.85) -- (3.85,4) node[pos = 1, above, xshift = 0.1cm] {\scriptsize$f_4$};
			\draw[xkcd, thick, c2] (-4,-4) -- (4,4) node [pos = 1, right] {$f$};
			\node[c1] (flabel) at (-3,2.5) {$f_n(x)=x+\frac{1}{n}$};
			\draw[xkcd, c1, dashed, ->] (flabel) to[bend left] (-0.4, 0.8);
		\end{tikzpicture}};
	
		\node[text width = 10cm] (label 1) at (7,16) {\raggedleft\LARGE\baselineskip=28pt Θεωρούμε τις εξής ευθείες:
			\[f_n(x)=x+\frac{1}{n}.\]
			Καθώς το $n$ μεγαλώνει και πάει προς το άπειρο, αυτές οι ευθείες προσεγγίζουν όλο και περισσότερο την ευθεία $f(x)=x$, όπως φαίνεται και δίπλα\ldots\par};
	
		\node[scale = 1.1] (first attempt) at (7,6) {\begin{tikzpicture}
			\draw[xkcd, thick, ->] (-1,0) -- (7,0) node[pos=1, below]{$x$};
			\draw[xkcd, thick, ->] (0,-1) -- (0,7) node[pos=1, left]{$y$};
			\node[point, c1] (zero) at (0,0) {};
			\node[point, c1] (one) at (6,6) {};
			% f1
			\draw[xkcd, c1, thick] (zero) -- (one) node[pos = 0.5, above, rotate = 45] {\scriptsize$f_1$};
			% f2
			\draw[xkcd, c1, thick, domain = {0}:{6}, opacity = 0.8] plot (\x,{pow(\x,2) / 6});
			\node[above, c1, opacity = 0.8, rotate = 45] (f2 label) at (3.5,2.0) {\scriptsize$f_2$};
			% f3
			\draw[xkcd, c1, thick, domain = {0}:{6}, opacity = 0.6] plot (\x,{pow(\x,3) / 36});
			\node[above, c1, opacity = 0.6, rotate = 48] (f3 label) at (4,1.75) {\scriptsize$f_3$};
			% f4
			\draw[xkcd, c1, thick, domain = {0}:{6}, opacity = 0.4] plot (\x,{pow(\x,4) / 216});
			\node[above, c1, opacity = 0.4, rotate = 52] (f4 label) at (4.42,1.65) {\scriptsize$f_4$};
			% f5
			\draw[xkcd, c1, thick, domain = {0}:{6}, opacity = 0.2] plot (\x,{pow(\x,5) / pow(6,4)});
			\node[above, c1, opacity = 0.2, rotate = 55] (f5 label) at (4.72,1.55) {\tiny$f_5$};
			\foreach \i in {6,...,20} {
				\pgfmathsetmacro{\opacity}{0.14 - (\i - 6) / 140}
				\draw[xkcd, c1, domain = {0}:{1}, opacity = \opacity, scale = 6] plot (\x,{pow(\x, {\i})});
			}
		\end{tikzpicture}};
	
		\node[text width = 10cm] (label 2) at (-7,6) {\raggedright\LARGE\baselineskip=28pt Ας πάρουμε τώρα τις συναρτήσεις που βλέπετε δίπλα. Αυτές δεν είναι άλλες παρά οι:
			\[f_n=x^n.\]
			Πού θα συγκλίνουν αυτές καθώς το $n$ πάει στο άπειρο;\par};
	
		\node[scale = 1.1] (second attempt) at (-7,-4) {\begin{tikzpicture}
			\draw[xkcd, thick, ->] (-1,0) -- (7,0) node[pos=1, below]{$x$};
			\draw[xkcd, thick, ->] (0,-1) -- (0,7) node[pos=1, left]{$y$};
			\node[point, c2] (zero) at (0,0) {};
			\node[point, c2] (one) at (6,6) {};
			\begin{pgfonlayer}{bg}
			% f1
			\draw[xkcd, c1, thick] (zero) -- (one) node[pos = 0.5, above, rotate = 45] {\scriptsize$f_1$};
			% f2
			\draw[xkcd, c1, thick, domain = {0}:{6}, opacity = 0.8] plot (\x,{pow(\x,2) / 6});
			\node[above, c1, opacity = 0.8, rotate = 45] (f2 label) at (3.5,2.0) {\scriptsize$f_2$};
			% f3
			\draw[xkcd, c1, thick, domain = {0}:{6}, opacity = 0.6] plot (\x,{pow(\x,3) / 36});
			\node[above, c1, opacity = 0.6, rotate = 48] (f3 label) at (4,1.75) {\scriptsize$f_3$};
			% f4
			\draw[xkcd, c1, thick, domain = {0}:{6}, opacity = 0.4] plot (\x,{pow(\x,4) / 216});
			\node[above, c1, opacity = 0.4, rotate = 52] (f4 label) at (4.42,1.65) {\scriptsize$f_4$};
			% f5
			\draw[xkcd, c1, thick, domain = {0}:{6}, opacity = 0.2] plot (\x,{pow(\x,5) / pow(6,4)});
			\node[above, c1, opacity = 0.2, rotate = 55] (f5 label) at (4.72,1.55) {\tiny$f_5$};
			\foreach \i in {6,...,20} {
				\pgfmathsetmacro{\opacity}{0.14 - (\i - 6) / 140}
				\draw[xkcd, c1, domain = {0}:{1}, opacity = \opacity, scale = 6] plot (\x,{pow(\x, {\i})});
			}
			\end{pgfonlayer}
			\draw[xkcd, ultra thick, c2] (0,0) -- (6,0) node[pos=0.5, below] {$f$};
			\node[empty point, draw = c2] (discontinuity) at (6,0) {};
			\draw[xkcd, dashed, c2] (discontinuity) -- (one);
		\end{tikzpicture}};
	
		\node[text width = 10cm] (label 3) at (7,-4) {\raggedleft\LARGE\baselineskip=28pt Με λίγο κόπο και λίγη προσπάθεια, μπορείτε να δείτε ότι συγκλίνουν στη συνάρτηση:
			\[f(x)=\begin{cases}
				0 & x\in[0,1)\\
				1 & x = 1
			\end{cases}\]\par};
		
		\node[text width = 15cm] (label 4) at (0,-15) {\centering\LARGE\baselineskip=28pt Τι διαφορετικό έχουν, άραγε, οι δύο παραπάνω περιπτώσεις;\par};
		
%		\node[scale = 1.1] (problem) at (7,-15) {\begin{tikzpicture}
%			
%		\end{tikzpicture}};
	
		\node (tbc) at (0,-22) {\Huge \textbf{Συνεχίζεται\ldots}};
	\end{tikzpicture}
\end{document}